\documentclass[11pt,a4paper]{article}

\usepackage{amsmath, amsthm, amssymb}
\usepackage{hyperref}
\usepackage{geometry}
\geometry{margin=2.5cm}

\newtheorem{proposition}{Proposition}
\newtheorem{lemma}{Lemma}
\newtheorem{remark}{Remark}
\newtheorem{definition}{Definition}

\title{Spectral Admissibility in the $Q_8$ Cayley Graph:\\
A Minimal Toy Model for Bounded-Flux Mode Selection}
\author{Working note -- Cosmochrony programme}
\date{\today}

\begin{document}
  \maketitle

  \begin{abstract}
    We establish the first analytically tractable instance of the spectral admissibility
    programme outlined in the Cosmochrony framework.
    On the undirected Cayley graph of the quaternion group $Q_8$, with generators
    $S = \{\pm i, \pm j, \pm k\}$, the local bounded-flux constraint $|\partial_t \chi_v| \leq c_\chi$
    induces a hierarchy of maximal admissible effective amplitudes across irreducible
    spectral sectors.
    In the linearised dispersion regime, the non-abelian irreducible sector
    admits a maximal effective amplitude larger than the nontrivial abelian sectors
    by a factor of $\sqrt{4/3} \approx 1.155$.
    This constitutes a minimal model demonstration that bounded substrate capacity
    may structurally favour non-abelian relational modes, independently of any dynamical
    postulate.
  \end{abstract}

  \section{Setting and Motivation}

    The Cosmochrony framework identifies the physical spectrum with the residue of two
    successive constraints acting on the relational operator $\mathcal{L}_\chi$:
    \begin{enumerate}
      \item \textbf{Topological discretisation}: the fibre structure of the non-injective
      projection $\Pi$ organises modes into discrete winding sectors.
      \item \textbf{Bounded-flux filtering}: the saturation bound
      $|\partial_t \chi| \leq c_\chi$ eliminates modes whose local evolution would
      exceed the substrate's relational capacity.
    \end{enumerate}
    The goal of this note is to make step~2 explicit and calculable on the simplest
    non-trivial finite graph encoding a quaternionic algebraic structure,
    as a minimal discrete proxy for $\mathrm{SU}(2)$-type sectors.

  \section{The $Q_8$ Cayley Graph}

    Let $Q_8 = \{{\pm 1, \pm i, \pm j, \pm k}\}$ denote the quaternion group of order~8.
    We define the undirected Cayley graph $\Gamma(Q_8, S)$ with generating set
    \begin{equation}
      S = \{\pm i, \pm j, \pm k\}.
    \end{equation}
    Since $S$ is closed under inversion ($s^{-1} \in S$ for all $s \in S$) and does
    not contain the identity, $\Gamma(Q_8, S)$ is a well-defined, connected, undirected,
    $6$-regular graph on $8$ vertices.

    \begin{remark}
      The choice $S = \{i, j, k\}$ (without inverses) would yield a directed graph with an
      asymmetric adjacency operator.
      For a self-adjoint Laplacian, the full symmetric set $S = \{\pm i, \pm j, \pm k\}$ is
      required.
      The resulting graph is $6$-regular, not $3$-regular.
    \end{remark}

    \subsection{Adjacency Operator and Graph Laplacian}

      The adjacency operator $A$ of $\Gamma(Q_8, S)$ acts on functions
      $f \colon Q_8 \to \mathbb{C}$ by
      \begin{equation}
        (Af)(g) = \sum_{s \in S} f(gs) = \sum_{s \in \{\pm i, \pm j, \pm k\}} f(gs).
      \end{equation}
      The combinatorial graph Laplacian is
      \begin{equation}
        \mathcal{L} = dI - A, \qquad d = |S| = 6,
      \end{equation}
      where $I$ is the identity on $\ell^2(Q_8)$.
      The eigenvalues of $\mathcal{L}$ are $\lambda = 6 - \mu$, where $\mu$ ranges over
      the eigenvalues of $A$.

  \section{Spectral Decomposition via Representation Theory}

    Since $Q_8$ is a finite group, $\ell^2(Q_8) \cong \mathbb{C}[Q_8]$ decomposes
    as a $Q_8$-module into irreducible representations.
    By the Peter--Weyl theorem for finite groups:
    \begin{equation}
      \mathbb{C}[Q_8] \cong \bigoplus_{\rho \in \widehat{Q_8}} V_\rho^{\oplus \dim \rho},
    \end{equation}
    where the sum runs over the five irreducible representations of $Q_8$.

    \subsection{Irreducible Representations of $Q_8$}

      $Q_8$ has five irreducible representations:
      \begin{center}
        \begin{tabular}{lcccc}
          \hline
          Label & Dimension & $\rho(i)$ & $\rho(j)$ & $\rho(k)$ \\
          \hline
          $\rho_0$ (trivial) & 1 & $1$ & $1$ & $1$ \\
          $\rho_1$ & 1 & $1$ & $-1$ & $-1$ \\
          $\rho_2$ & 1 & $-1$ & $1$ & $-1$ \\
          $\rho_3$ & 1 & $-1$ & $-1$ & $1$ \\
          $\rho_4$ (quaternionic) & 2 & $\begin{pmatrix}i&0\\0&-i\end{pmatrix}$
          & $\begin{pmatrix}0&1\\-1&0\end{pmatrix}$
          & $\begin{pmatrix}0&i\\i&0\end{pmatrix}$ \\
          \hline
        \end{tabular}
      \end{center}
      The dimension check: $1^2 + 1^2 + 1^2 + 1^2 + 2^2 = 8 = |Q_8|$. \checkmark

    \subsection{Eigenvalues of the Adjacency Operator}

      For any irreducible representation $\rho$, the adjacency operator acts as the
      scalar $\mu_\rho \cdot I_{V_\rho}$ on $V_\rho$.
      Since $S$ is a union of conjugacy classes of $Q_8$, the corresponding element
      of the group algebra is central, and Schur's lemma gives
      \begin{equation}
        \mu_\rho = \frac{1}{\dim \rho} \sum_{s \in S} \chi_\rho(s),
      \end{equation}
      where $\chi_\rho$ is the character of $\rho$.

      \paragraph{1D sectors.}
        For the trivial representation $\rho_0$: $\chi_{\rho_0}(s) = 1$ for all $s \in S$,
        so $\sum_{s \in S} \chi_{\rho_0}(s) = 6$ and
      \begin{equation}
        \mu_{\rho_0} = \frac{1}{1} \cdot 6 = 6, \qquad \lambda_{\rho_0} = 6 - 6 = 0.
        \end{equation}
        For $\rho_1$: $\chi_{\rho_1}(\pm i) = 1$ and $\chi_{\rho_1}(\pm j) = \chi_{\rho_1}(\pm k) = -1$,
        so $\sum_{s \in S} \chi_{\rho_1}(s) = 2(1) + 4(-1) = -2$, giving
        $\mu_{\rho_1} = -2$ and $\lambda_{\rho_1} = 8$.
        By the symmetry of $Q_8$ under permutations of $\{i,j,k\}$,
        $\mu_{\rho_2} = \mu_{\rho_3} = -2$ and $\lambda_{\rho_2} = \lambda_{\rho_3} = 8$.

      \paragraph{2D sector.}
        For the quaternionic representation $\rho_4$, we compute the character at each
        generator.
        The matrices $\rho_4(i)$, $\rho_4(-i)$, $\rho_4(j)$, $\rho_4(-j)$, $\rho_4(k)$,
        $\rho_4(-k)$ each have trace zero: for instance,
        $\mathrm{tr}(\rho_4(i)) = \mathrm{tr}\bigl(\begin{smallmatrix}i&0\\0&-i\end{smallmatrix}\bigr) = 0$,
        and $\mathrm{tr}(\rho_4(-i)) = 0$; similarly for $j$ and $k$.
        Thus $\sum_{s \in S} \chi_{\rho_4}(s) = 0$ and
      \begin{equation}
        \mu_{\rho_4} = \frac{1}{2} \cdot 0 = 0, \qquad \lambda_{\rho_4} = 6 - 0 = 6.
        \end{equation}

  \subsection{Summary of the Spectrum}

    \begin{center}
      \begin{tabular}{lccc}
        \hline
        Sector & $\lambda_n$ & Multiplicity & Type \\
        \hline
        $\rho_0$ & $0$ & $1$ & Trivial (zero mode) \\
        $\rho_1, \rho_2, \rho_3$ & $8$ & $3$ & Nontrivial abelian \\
        $\rho_4$ & $6$ & $4$ & Non-abelian (quaternionic) \\
        \hline
      \end{tabular}
    \end{center}
    Total dimension: $1 + 3 + 4 = 8$. \checkmark

  \section{Bounded-Flux Admissibility Condition}

  \subsection{Setup}

    Let the substrate field $\chi \colon Q_8 \times \mathbb{R} \to \mathbb{R}$ evolve on the
    graph, subject to the local bounded-flux constraint
    \begin{equation}
      \label{eq:flux-bound}
      |\partial_t \chi_v(t)| \leq c_\chi \qquad \forall v \in Q_8, \; \forall t.
    \end{equation}
    Consider a single-mode excitation
    \begin{equation}
      \chi_v(t) = a_n \psi_n(v) \cos(\omega_n t),
    \end{equation}
    where $\psi_n$ is an eigenmode of $\mathcal{L}$ with eigenvalue $\lambda_n$,
    normalised in $\ell^2(Q_8)$, and $a_n > 0$ is the mode amplitude.
    The time derivative at vertex $v$ is
    \begin{equation}
      \partial_t \chi_v(t) = -a_n \omega_n \psi_n(v) \sin(\omega_n t).
    \end{equation}
    The constraint~\eqref{eq:flux-bound} requires
    \begin{equation}
      |\omega_n| \cdot a_n \cdot |\psi_n(v)| \leq c_\chi \qquad \forall v \in Q_8.
    \end{equation}
    The binding constraint is at the vertex where $|\psi_n(v)|$ is maximal:
    \begin{equation}
      |\omega_n| \cdot a_n \cdot \|\psi_n\|_\infty \leq c_\chi.
    \end{equation}

  \subsection{Effective Amplitude and Admissibility Envelope}

    \begin{definition}
      The \emph{effective amplitude} of a mode $(\psi_n, a_n)$ is
      \begin{equation}
        A_n := a_n \|\psi_n\|_\infty.
      \end{equation}
      A mode is \emph{dynamically admissible} (in the linearised regime) if
      \begin{equation}
        \label{eq:admissibility}
        |\omega_n| \cdot A_n \leq c_\chi.
      \end{equation}
    \end{definition}

    In the linearised dispersion approximation, $\omega_n^2 \approx \lambda_n$, so the
    maximal admissible effective amplitude is
    \begin{equation}
      \label{eq:envelope}
      A_n^{\max} = \frac{c_\chi}{\sqrt{\lambda_n}}.
    \end{equation}
    This defines the \emph{admissibility envelope} in the $(\lambda_n, A_n)$ plane.

    \begin{remark}
      The linearised dispersion $\omega_n^2 = \lambda_n$ is valid for $A_n \ll c_\chi / \sqrt{\lambda_n}$.
      The non-linear (Born-Infeld) correction modifies this relation near saturation and
      requires a modal reduction of the DBI action; it is deferred to a separate calculation.
      The envelope~\eqref{eq:envelope} should therefore be understood as the admissibility
      boundary in the weakly non-linear regime.
    \end{remark}

  \section{Main Result}

  \begin{proposition}[Spectral admissibility hierarchy on $Q_8$]
    \label{prop:main}
    In the undirected Cayley graph $\Gamma(Q_8, \{\pm i, \pm j, \pm k\})$,
    under the local bounded-flux constraint~\eqref{eq:flux-bound} and linearised
    dispersion, the maximal admissible effective amplitude satisfies
    \begin{equation}
      A^{\max}_{\rho_4} = \frac{c_\chi}{\sqrt{6}} > \frac{c_\chi}{\sqrt{8}}
      = A^{\max}_{\rho_k}, \quad k = 1, 2, 3.
    \end{equation}
    The non-abelian sector $\rho_4$ occupies a strictly larger admissible region in the
    space of effective mode amplitudes than any nontrivial abelian sector, by a factor
    \begin{equation}
      \frac{A^{\max}_{\rho_4}}{A^{\max}_{\rho_k}} = \sqrt{\frac{8}{6}} = \sqrt{\frac{4}{3}} \approx 1.155.
    \end{equation}
  \end{proposition}

  \begin{proof}
    Direct substitution of $\lambda_{\rho_4} = 6$ and $\lambda_{\rho_k} = 8$
    into the admissibility envelope~\eqref{eq:envelope}.
  \end{proof}

  \begin{remark}
    The trivial sector $\rho_0$ has $\lambda = 0$ and thus formally unbounded $A^{\max}$.
    It corresponds to the uniform zero mode of the connected graph.
    In the present analysis it is factored out, since we are interested only in
    non-trivial excitations above the connected background.
  \end{remark}

  \section{Interpretation}

  Proposition~\ref{prop:main} establishes that, at fixed substrate capacity $c_\chi$,
  the non-abelian irreducible sector of the quaternion Cayley graph sustains a larger
  local effective amplitude before violating the bounded-flux constraint.
  In the language of the Cosmochrony programme:

  \begin{itemize}
    \item The admissibility envelope $A_n^{\max} = c_\chi / \sqrt{\lambda_n}$ defines a
    \emph{geometry of exclusion} in the space of relational modes: modes above the
    envelope are dynamically inadmissible.
    \item The non-abelian sector $\rho_4$ lies strictly below this envelope for a larger
    range of effective amplitudes than the nontrivial abelian sectors
    $\rho_1, \rho_2, \rho_3$: the bounded-flux constraint is less restrictive for
    non-abelian modes, enlarging their admissible window by a factor of $\sqrt{4/3}$.
    \item This provides a minimal admissibility advantage for non-abelian relational
    modes, without any dynamical postulate beyond the local flux bound.
  \end{itemize}

  \begin{remark}[Real field and complex eigenmodes]
    The substrate field $\chi$ is taken real-valued.
    The natural eigenmodes of the 2D sector are matrix coefficients of $\rho_4$,
    which are complex-valued.
    One may either work in $\ell^2(Q_8, \mathbb{C})$ and take real and imaginary parts to
    obtain real eigenmodes, or directly choose a real basis for each eigenspace.
    The admissibility bound on $A_n$ is basis-independent (it depends on
    $\|\psi_n\|_\infty$, which may vary within a degenerate eigenspace); this basis
    dependence is addressed in Step~1 of Section~\ref{sec:next-steps}.
  \end{remark}

  This result does not yet establish:
  \begin{itemize}
    \item that the non-abelian sector will dominate in a full dynamical simulation;
    \item the origin of the eigenvalue gap $\lambda_{\rho_4} = 6 < 8 = \lambda_{\rho_k}$
    from the topological structure of $\Pi$ alone;
    \item the Born-Infeld correction to the admissibility envelope near saturation.
  \end{itemize}
  These remain open and constitute the next steps of the programme.

  \section{Next Steps}

  The natural continuation of this calculation proceeds as follows.

  \paragraph{Step 1: explicit eigenmodes and $\|\psi_n\|_\infty$.}
    \label{sec:next-steps}
    Choose an explicit $\ell^2$-normalised basis in each eigenspace using the matrix
    coefficients of the representations.
    For each basis vector, compute $\|\psi_n\|_\infty$ explicitly.
    In degenerate eigenspaces (multiplicity $> 1$, in particular the 2D sector with
    multiplicity~4), determine the basis dependence of the admissibility bound
    $a_n^{\max} = c_\chi / (\sqrt{\lambda_n}\,\|\psi_n\|_\infty)$ and identify the
    basis minimising or maximising $\|\psi_n\|_\infty$ within the eigenspace.
    Expected values: $\|\psi_n\|_\infty = 1/\sqrt{8}$ for 1D sectors and
    $\|\psi_n\|_\infty \leq 1/\sqrt{2}$ for the 2D sector (upper bound from
    $L^2$-normalisation).

  \paragraph{Step 2: modal DBI reduction.}
    For a single mode $\chi_v(t) = q(t)\psi_n(v)$, substitute into the discrete DBI
    Lagrangian
    \begin{equation}
      L = -c_\chi^2 \sum_v \sqrt{1 - \dot{\chi}_v^2/c_\chi^2}
      - \tfrac{1}{2} \sum_{v,w} \chi_v \mathcal{L}_{vw} \chi_w,
    \end{equation}
    and derive the effective equation for $q(t)$.
    Extract the amplitude-dependent frequency $\omega_n(a_n)$ via harmonic balance.
    This yields the corrected admissibility envelope replacing~\eqref{eq:envelope}.

  \paragraph{Step 3: dynamical simulation.}
    Simulate the full $Q_8$ substrate with the DBI constraint imposed locally at
    each vertex.
    Track the modal decomposition $A_n(t) = a_n(t)\|\psi_n\|_\infty$ over time.
    Test whether modes initialised above the envelope~\eqref{eq:envelope} decay,
    and whether modes in the non-abelian sector survive longer than those in
    nontrivial abelian sectors at comparable initial $A_n$.

  \paragraph{Step 4: extension to larger graphs.}
    Replace $Q_8$ by the Cayley graph of a group approximating $S^3$ more faithfully
    at large $N$ (e.g.\ the binary icosahedral group $2I$ of order~120, or
    $\mathrm{SL}(2, \mathbb{F}_p)$ for small primes $p$).
    Test whether the spectral admissibility hierarchy persists under this refinement.
    A particularly natural continuation is provided by Cayley graphs in Ramanujan
    families -- such as the Lubotzky--Phillips--Sarnak construction on
    $\mathrm{SL}(2, \mathbb{F}_p)$ -- where the spectral gap is asymptotically optimal.
    This would allow the admissibility hierarchy to be tested in graph sequences with
    extremal expansion properties, directly connecting bounded-flux filtering to the
    theory of expander graphs.

\end{document}
